\textbf{Proof.}
Fix \(0<\epsilon<1/2\) and write \(L_n=\log(en)\).  We first globalize the two lower-bound mechanisms.

By \(\mathrm{lem:paired-sign-small-alphabet-floor}\), for \(2\le d\le n^2\),
\[
 \mathfrak R_{n,d,\epsilon}\ge c_0\frac{\sqrt d}{n}
\]
for a universal \(c_0>0\).  If \(n\ge2\) and \(d>n^2\), apply \(\mathrm{lem:alphabet-monotonicity}\) with \(m=n^2\) and then apply the same paired-sign lemma at alphabet size \(m\).  This gives
\[
 \mathfrak R_{n,d,\epsilon}
 \ge\mathfrak R_{n,n^2,\epsilon}\ge c_0.
\]
For \(n=1\), the parametric floor in \(\mathrm{lem:l1-lower-transfer}\) supplies a positive constant.  After decreasing the constant, these cases combine as
\[
 \mathfrak R_{n,d,\epsilon}
 \ge c_{1,\epsilon}\min\left\{1,\frac{\sqrt d}{n}\right\}.
\]

Next let \(c_{2,\epsilon}>0\) and \(\rho_0>0\) be the constants in \(\mathrm{lem:l1-lower-transfer}\).  On \(2\le d\le\rho_0nL_n\), that lemma gives
\[
 \mathfrak R_{n,d,\epsilon}
 \ge c_{2,\epsilon}\frac d{nL_n}.
\]
Choose \(N_0\) so that \(\rho_0nL_n\ge8\) for \(n\ge N_0\).  If \(n\ge N_0\) and \(d>\rho_0nL_n\), set
\[
 m=\left\lfloor\frac{\rho_0nL_n}{2}\right\rfloor.
\]
Then \(2\le m\le d\), \(m\le\rho_0nL_n\), and \(m\ge\rho_0nL_n/4\).  Alphabet monotonicity and the in-range lower bound give
\[
 \mathfrak R_{n,d,\epsilon}
 \ge\mathfrak R_{n,m,\epsilon}
 \ge c_{2,\epsilon}\frac m{nL_n}
 \ge\frac{c_{2,\epsilon}\rho_0}{4}.
\]
For the finitely many \(n<N_0\), the parametric floor \(1/(1024n)\) is at least \(1/(1024N_0)\).  Decreasing the constant once more yields, for all \(n\ge1,d\ge2\),
\[
 \mathfrak R_{n,d,\epsilon}
 \ge c_{3,\epsilon}\min\left\{1,\frac d{nL_n}\right\}.
\]

For nonnegative \(u,v\),
\[
 \max\{\min(1,u),\min(1,v)\}
 \ge\frac12\{\min(1,u)+\min(1,v)\}
 \ge\frac12\min(1,u+v).
\]
Applying this with \(u=\sqrt d/n\) and \(v=d/(nL_n)\), and taking the smaller lower constant, proves
\[
 \mathfrak R_{n,d,\epsilon}
 \ge c_\epsilon\min\left\{1,
   \frac{\sqrt d}{n}+\frac d{nL_n}
 \right\}.
\]

For the upper bound, the definition of minimax risk permits evaluation at the single data-driven estimator \(\widehat V_n^{\mathrm{DAT}}\).  Lemma \(\mathrm{lem:empirical-cellwise-ratio-upper}\) then gives
\[
 \mathfrak R_{n,d,\epsilon}
 \le\sup_{P\in\mathcal P_{d,\epsilon}}
 E_P\left[
   \{\widehat V_n^{\mathrm{DAT}}-V^\star(P)\}^2
 \right]
 \le\min\left\{1,
  \left(\frac12+\frac6\epsilon\right)\frac dn
 \right\}.
\]
Thus one may take \(C_\epsilon=1/2+6/\epsilon\), enlarged if necessary to satisfy \(c_\epsilon\le C_\epsilon\).

The ratio of the nonsaturated lower terms is
\[
 \frac{\sqrt d/n}{d/(nL_n)}=\frac{L_n}{\sqrt d},
\]
so they balance exactly at order \(d\asymp L_n^2\), with the claimed dominance on either side.

Consider now sequences indexed by \(n\to\infty\).  If the minimax risk tends to zero, the second global lower bound forces
\[
 \frac d{nL_n}\longrightarrow0,
\]
which is equivalent to \(d=o(n\log n)\).  Conversely, \(d=o(n)\) makes the displayed upper bound tend to zero.

Finally, if \(d=O(1)\), the upper bound is \(O(n^{-1})\), while the parametric floor in \(\mathrm{lem:l1-lower-transfer}\) is \(\Omega(n^{-1})\).  Hence the risk is \(\Theta(n^{-1})\).  Conversely, suppose the risk is \(O(n^{-1})\).  If \(d>n^2\) along a subsequence, the first globalized lower bound is bounded away from zero, a contradiction.  Otherwise \(d\le n^2\) eventually, and the paired-sign lower bound gives \(c_0\sqrt d/n=O(n^{-1})\), so \(d=O(1)\).  This proves the iff statement and also shows precisely why the displayed bracket leaves the intermediate consistency region and matching achievability unresolved.